\documentclass[11pt]{article}
\usepackage[margin=1in]{geometry}
\usepackage{amsmath,amssymb}
\usepackage{graphicx}
\usepackage{booktabs}
\usepackage{hyperref}
\usepackage{microtype}
\usepackage{enumitem}
\usepackage{float}
\usepackage{parskip}

\hypersetup{colorlinks=true,linkcolor=blue,urlcolor=blue,citecolor=blue}
\graphicspath{{figures/}}
\setlength{\parskip}{5pt}
\setlength{\parindent}{0pt}

\title{\textbf{Scaling Laws for Language Models Trained on SVG Code}\\[4pt]
\large CS-GY 6923 Optional Project --- Spring 2026}
\author{Tanmay Sahu \quad \texttt{ts5888@nyu.edu}\\
NYU Tandon School of Engineering\\[2pt]
\small \url{https://github.com/7dracoder/ML-Project}}
\date{May 1, 2026}

\begin{document}
\maketitle
\thispagestyle{plain}

%───────────────────────────────────────────────────────────────
\section{Introduction}

Scaling laws describe how model performance improves with model size, data, and
compute. Kaplan et al.\ (2020) established power-law relationships for natural
language, and Hoffmann et al.\ (2022) refined these with compute-optimal recipes.
Scaling behavior in structured, non-linguistic domains remains underexplored.

SVG (Scalable Vector Graphics) is XML-based text encoding visual content.
Unlike natural language, SVG has strict syntactic rules, hierarchical structure,
and outputs that can be instantly rendered --- making evaluation both quantitative
and qualitative. We train five decoder-only transformers (1.4M--88.6M params) on
115M tokens of SVG data, fit a power-law scaling curve, investigate $\mu$P for
learning rate transfer, and generate and evaluate SVG samples from our best model.

%───────────────────────────────────────────────────────────────
\section{Data}

\subsection{Datasets and Preprocessing}

We use five HuggingFace datasets: \texttt{svg-icons-simple} (80K, primary),
\texttt{svg-emoji-simple} (4K), \texttt{svg-fonts-simple} (100K subsampled),
\texttt{svg-stack-simple} (120K subsampled), and \texttt{svgen-500k} (attempted
but unavailable on Hub). After MD5 deduplication: 300,045 unique SVGs; after
token-length filtering ($\leq$1024 tokens): 223,084.

Each SVG is cleaned by: stripping comments/declarations/DOCTYPE, collapsing
whitespace, rounding coordinates to 1 decimal place (e.g.\ $45.2387 \to 45.2$),
validating XML with \texttt{lxml}, and removing SVGs under 50 chars or over 1024 tokens.

\subsection{Tokenization}

We train a BPE tokenizer (HuggingFace \texttt{tokenizers}, vocab=4096).
BPE learns SVG-specific patterns as tokens --- e.g.\ \texttt{<path d="M} becomes
one token --- helping the model learn structure rather than character-by-character
syntax. Vocab 4096 balances compression (2.4 chars/token) with embedding table
size (38\% of Tiny's parameters at this vocab; 72\% at 8192 --- too dominant).

\subsection{Statistics}

\begin{table}[H]
\centering\small
\begin{tabular}{lr}
\toprule
Metric & Value \\
\midrule
Vocabulary size & 4,096 \\
Train tokens & 112,654,174 \\
Val / Test tokens & 1,139,210 / 1,165,205 \\
Mean / Median seq length & 515 / 481 tokens \\
\bottomrule
\end{tabular}
\caption{Dataset statistics. Split is 98/1/1 by file to prevent leakage.}
\end{table}

\begin{figure}[H]
\centering
\includegraphics[width=0.85\linewidth]{seq_lengths.png}
\caption{Sequence length distribution and CDF.}
\end{figure}

\begin{figure}[H]
\centering
\includegraphics[width=0.85\linewidth]{svg_examples.png}
\caption{Rendered SVG examples at P10 (simple), P50 (medium), P90 (complex) complexity.}
\end{figure}

%───────────────────────────────────────────────────────────────
\section{Methods}

\subsection{Model Architecture}

Decoder-only GPT with pre-norm, Flash Attention (\texttt{F.scaled\_dot\_product\_attention}),
GELU, weight tying (SP), and residual scaling ($\sigma = 0.02/\sqrt{2L}$).

\begin{table}[H]
\centering\small
\begin{tabular}{lrrrrr}
\toprule
Model & Params & $d$ & $L$ & $H$ & $d_{ff}$ \\
\midrule
Tiny   & 1.38M  & 128 & 4  & 4  & 512  \\
Small  & 3.55M  & 192 & 6  & 6  & 768  \\
Medium & 12.4M  & 384 & 6  & 6  & 1536 \\
Large  & 33.9M  & 512 & 10 & 8  & 2048 \\
XL     & 88.6M  & 768 & 12 & 12 & 3072 \\
\bottomrule
\end{tabular}
\caption{Model configurations. All use context length 512 tokens.}
\end{table}

\subsection{Training Setup}

AdamW ($\beta_1=0.9$, $\beta_2=0.95$, wd=0.1), cosine LR with 10\% warmup,
batch 32,768 tokens, bfloat16, grad clip 1.0, dropout 0.0 (scaling study) /
0.1 (extended training). All models trained for exactly 1 epoch for the scaling
comparison.

\textbf{LR strategy (Part 2):} The Tiny LR sweep (7 candidates, $3\times10^{-4}$ to
$8\times10^{-3}$) found $8\times10^{-3}$ optimal. Per the spec, we first apply this
same LR to all sizes (``fixed LR''), showing the expected degradation at Large/XL.
We then use the standard sqrt-scaling heuristic
$\text{lr}(N) = \text{lr}_{\text{tiny}} \cdot \sqrt{N_{\text{tiny}}/N}$
to produce a proper SP scaling curve. This directly motivates $\mu$P in Part 3.

\subsection{$\mu$P Implementation}

We use the \texttt{mup} package with: (1) attention scaling $1/d_{\text{head}}$
instead of $1/\sqrt{d_{\text{head}}}$; (2) \texttt{MuReadout} output layer (no
weight tying --- required by $\mu$P); (3) \texttt{MuAdam} optimizer; (4) separate
Q/K/V linear projections for proper shape tracking via \texttt{set\_base\_shapes}.
Separate LR sweep on Tiny $\mu$P; optimal LR ($8\times10^{-3}$) transferred
to all larger $\mu$P models without retuning.

\subsection{Evaluation Metrics}

Validation/test loss and perplexity; XML validity rate (\texttt{lxml}); SVG render
rate (CairoSVG); structural validity (SVG root, closed tags, valid attributes);
wall-clock time, GPU memory, tokens/second.

%───────────────────────────────────────────────────────────────
\section{Results}

\subsection{Scaling Plot and Power Law Fit}

\begin{table}[H]
\centering\small
\begin{tabular}{lrrr}
\toprule
Model & Params & Fixed LR Val & Scaled LR Val \\
\midrule
Tiny   & 1.38M  & 0.5898 & 0.5834 \\
Small  & 3.55M  & 0.5398 & 0.5348 \\
Medium & 12.4M  & 0.5407 & 0.5266 \\
Large  & 33.9M  & 0.6738 & 0.5159 \\
XL     & 88.6M  & 1.4826 & 0.5017 \\
\bottomrule
\end{tabular}
\caption{Validation loss after 1 epoch. Fixed LR: same $8\times10^{-3}$ for all
sizes (per spec). Scaled LR: sqrt-heuristic. Fixed LR degrades at Large/XL ---
the expected SP failure mode that motivates $\mu$P.}
\end{table}

\begin{equation}
L = 7190.98 \cdot N^{-0.8107} + 0.5061, \quad R^2 = 0.9559
\end{equation}

\begin{figure}[H]
\centering
\includegraphics[width=0.85\linewidth]{scaling_final.png}
\caption{Left: SP and $\mu$P scaling curves with power-law fits and 10$\times$ extrapolation.
Right: LR sweep comparison (SP vs.\ $\mu$P) on Tiny model.}
\end{figure}

\subsection{Fixed LR vs.\ $\mu$P and LR Sweep}

\begin{figure}[H]
\centering
\includegraphics[width=0.85\linewidth]{lr_sweep.png}
\caption{LR sweep on Tiny. Best LR = $8\times10^{-3}$ (val=0.5888).}
\end{figure}

\begin{table}[H]
\centering\small
\begin{tabular}{lrr}
\toprule
Model & SP Val & $\mu$P Val \\
\midrule
Tiny   & 0.5834 & 2.4609 \\
Small  & 0.5348 & 2.7032 \\
Medium & 0.5266 & 3.9972 \\
Large  & 0.5159 & 4.0330 \\
XL     & 0.5017 & 4.1961 \\
\bottomrule
\end{tabular}
\caption{SP (scaled LR) vs.\ $\mu$P after 1 epoch. $\mu$P failed to achieve
competitive performance (see Discussion).}
\end{table}

\subsection{Scaling Law Extrapolation}

\begin{equation}
\hat{L}(885\text{M params}) = 0.4947, \quad 95\%\ \text{CI:}\ [0.4745,\ 0.5187]
\end{equation}
Bootstrap CI from 2000 resamples of the SP (scaled LR) power law fit.

\subsection{Training Curves}

\begin{figure}[H]
\centering
\includegraphics[width=0.85\linewidth]{extended_training.png}
\caption{XL model over 3 epochs (10,311 steps). Val loss: 0.65 $\to$ 0.45.}
\end{figure}

\subsection{Generated Samples and Evaluation}

Best model: XL, 3 epochs, dropout 0.1. Test perplexity: \textbf{1.57}.

\begin{table}[H]
\centering\small
\begin{tabular}{lr}
\toprule
Metric & Rate \\
\midrule
XML valid / SVG root / Closed tags & 54\% \\
Has viewBox / Valid attrs / Structural & 54\% \\
Renderable PNG (CairoSVG) & 54\% \\
\bottomrule
\end{tabular}
\caption{Generation quality on 100 samples (with SVG-closing post-processing).
All metrics coincide because XML parse failure cascades to all downstream checks.}
\end{table}

\begin{figure}[H]
\centering
\includegraphics[width=0.85\linewidth]{unconditional_grid.png}
\caption{Unconditional samples at $T=0.5$ (top), $0.8$ (middle), $1.0$ (bottom).}
\end{figure}

\begin{figure}[H]
\centering
\includegraphics[width=0.85\linewidth]{prefix_grid.png}
\caption{Prefix completions. Leftmost column: prefix. Columns 2--4: completions at $T \in \{0.5, 0.8, 1.0\}$.}
\end{figure}

%───────────────────────────────────────────────────────────────
\section{Discussion}

\subsection{Key Insights from Scaling Experiments}

Our $\alpha = 0.8107$ is steeper than both Kaplan et al.'s $\alpha \approx 0.076$
for natural language and Chinchilla's $\alpha \approx 0.34$. SVG's structured
syntax constrains the next-token distribution heavily --- given a partial SVG, valid
continuations are far fewer than in natural language. This lower entropy means
additional parameters help more per unit compute. The irreducible loss $c = 0.5061$
reflects inherent randomness in SVG data (exact coordinates, color choices) that
no model can predict.

\subsection{Learning Rate Scaling and $\mu$P}

Using the same LR for all model sizes under SP causes Large/XL to diverge ---
the effective LR for different layers scales differently with width. Our sqrt-scaled
LR heuristic resolves this. $\mu$P is designed to eliminate this heuristic entirely
by enabling zero-shot LR transfer.

Our $\mu$P implementation did not achieve competitive performance: Medium/Large/XL
models plateau at loss $\sim$4.0, indicating near-random predictions. The
\texttt{mup} package's \texttt{set\_base\_shapes} tracks layers by Python attribute
paths, and our custom architecture (\texttt{MuPAttention} with separate Q/K/V
inside \texttt{MuPBlock}) does not fully match the expected naming conventions.
\texttt{MuAdam} therefore does not apply correct per-layer LR multipliers for
wider models, effectively behaving as standard Adam with the wrong attention
scaling ($1/d_{\text{head}}$ instead of $1/\sqrt{d_{\text{head}}}$). A correct
implementation would require using \texttt{mup.MuLinear} layers throughout or
adapting the architecture to match \texttt{mup}'s expectations exactly.

\subsection{Design Decisions}

\textbf{Vocab 4096:} Balances compression with embedding table size.
\textbf{Context 512:} Covers median SVG (481 tokens); 1024 would quadruple attention memory.
\textbf{Pre-norm:} More stable than post-norm for deep models.
\textbf{No dropout in scaling study:} Ensures clean comparison across model sizes.
\textbf{SVG closing:} Generated SVGs that hit the token limit without closing tags
are automatically closed via post-processing, improving render rates.

\subsection{Limitations and Future Work}

115M tokens is far below Chinchilla-optimal for XL ($\sim$17.7B tokens). More data
would improve both scaling exponent and generation quality. Constrained decoding
enforcing XML validity would significantly improve the 54\% render rate. Increasing
context to 1024 tokens would cover 95\%+ of SVGs. A correct $\mu$P implementation
using \texttt{mup.MuLinear} layers would enable proper zero-shot LR transfer.

%───────────────────────────────────────────────────────────────
\section{Conclusion}

We trained decoder-only transformers on SVG code and found
$L = 7190.98 \cdot N^{-0.8107} + 0.5061$ ($R^2 = 0.9559$) across 1.4M--88.6M
parameters. The exponent $\alpha = 0.8107$ is steeper than natural language,
consistent with SVG's structured, lower-entropy nature. Our best model (XL, 3
epochs) achieves test perplexity 1.57 and 54\% render rate. $\mu$P did not improve
over SP due to implementation constraints with the \texttt{mup} library. The
extrapolated loss for a 10$\times$ larger model (885M params) is
$0.4947\ [0.4745, 0.5187]_{95\%}$.

%───────────────────────────────────────────────────────────────
\appendix
\section*{Appendix: Summary Figure}
\begin{figure}[H]
\centering
\includegraphics[width=0.9\linewidth]{scaling_final.png}
\caption{Scaling laws and LR sweep comparison.}
\end{figure}

\end{document}
